\documentclass[12pt,a4paper]{article}
\usepackage{fontspec}
\setmainfont{Times New Roman}
\usepackage{polyglossia}
\setdefaultlanguage{russian}
\usepackage{amsmath,amssymb}
\usepackage{geometry}
\usepackage{booktabs}
\usepackage{array}
\usepackage{float}
\geometry{left=25mm,right=15mm,top=20mm,bottom=20mm}
\setlength{\parindent}{1.25cm}
\setlength{\parskip}{0pt}

\begin{document}

\begin{titlepage}
\begin{center}
ФЕДЕРАЛЬНОЕ АГЕНТСТВО ПО ОБРАЗОВАНИЮ\\[3mm]
ИТМО\\[35mm]
\textbf{ЛАБОРАТОРНАЯ РАБОТА № 6}\\[5mm]
\textbf{ЧИСЛЕННОЕ РЕШЕНИЕ ОБЫКНОВЕННЫХ\\
ДИФФЕРЕНЦИАЛЬНЫХ УРАВНЕНИЙ}\\[35mm]
\end{center}

\begin{flushright}
Выполнил: студент группы P3215\\
Панченко А. Д.
\end{flushright}

\vfill
\begin{center}
Санкт-Петербург\\2026
\end{center}
\end{titlepage}

\section*{Цель работы}
Получить численное решение задачи Коши для обыкновенного дифференциального уравнения, реализовав методы Эйлера, Рунге--Кутты четвёртого порядка и Милна. Сравнить полученные значения с точным решением и оценить погрешность.

\section{Постановка задачи}
Вариант 10 предусматривает использование методов 1, 3 и 5:
\begin{enumerate}
    \item метод Эйлера;
    \item метод Рунге--Кутты четвёртого порядка;
    \item метод Милна.
\end{enumerate}

В программной реализации доступны несколько задач Коши. Для контрольного расчёта выбрана задача
\begin{equation}
    y'=-y, \qquad y(0)=1, \qquad x\in[0;2], \qquad h=0{,}1.
\end{equation}
Её точное решение имеет вид
\begin{equation}
    y(x)=e^{-x}.
\end{equation}

Сетка вычислений задаётся формулами
\begin{equation}
    x_i=x_0+ih, \qquad i=0,1,\ldots,n, \qquad n=\frac{x_n-x_0}{h}.
\end{equation}

\section{Теоретические сведения}
\subsection{Метод Эйлера}
Метод Эйлера является одношаговым методом первого порядка точности. Значение в следующем узле вычисляется по формуле
\begin{equation}
    y_{i+1}=y_i+h f(x_i,y_i).
\end{equation}
Для выбранной задачи $f(x,y)=-y$, поэтому рекуррентная формула принимает вид
\begin{equation}
    y_{i+1}=y_i(1-h).
\end{equation}

\subsection{Метод Рунге--Кутты четвёртого порядка}
Для каждого шага вычисляются четыре вспомогательных значения:
\begin{align}
 k_1&=h f(x_i,y_i),\\
 k_2&=h f\left(x_i+\frac h2,y_i+\frac{k_1}{2}\right),\\
 k_3&=h f\left(x_i+\frac h2,y_i+\frac{k_2}{2}\right),\\
 k_4&=h f(x_i+h,y_i+k_3).
\end{align}
Затем
\begin{equation}
    y_{i+1}=y_i+\frac{k_1+2k_2+2k_3+k_4}{6}.
\end{equation}
Метод имеет локальную погрешность $O(h^5)$ и глобальную погрешность $O(h^4)$.

\subsection{Метод Милна}
Метод Милна является многошаговым методом четвёртого порядка. Для начала работы ему требуются четыре начальные точки. Значения $y_1$, $y_2$, $y_3$ в программе вычисляются методом Рунге--Кутты четвёртого порядка.

На очередном шаге сначала вычисляется прогноз:
\begin{equation}
 y_{i+1}^{\text{пр}}=y_{i-3}+\frac{4h}{3}\left(2f_{i-3}-f_{i-2}+2f_{i-1}\right),
 \qquad f_j=f(x_j,y_j).
\end{equation}
После этого прогноз подставляется в корректирующую формулу:
\begin{equation}
 y_{i+1}=y_{i-2}+\frac h3\left(f_{i-2}+4f_{i-1}+f(x_{i+1},y_{i+1}^{\text{пр}})\right).
\end{equation}

\section{Оценка погрешности}
Для методов Эйлера и Рунге--Кутты используется правило Рунге:
\begin{equation}
 R=\frac{|y_h(x_n)-y_{h/2}(x_n)|}{2^p-1},
\end{equation}
где $p$ --- порядок метода. Дополнительно для всех методов вычисляется максимальная абсолютная погрешность по точному решению:
\begin{equation}
 \varepsilon_{\max}=\max_i|y_i-y(x_i)|.
\end{equation}

\section{Пример вычисления}
Для первого шага метода Эйлера при $x_0=0$, $y_0=1$, $h=0{,}1$:
\begin{equation}
 y_1=y_0+h(-y_0)=1-0{,}1=0{,}9.
\end{equation}

Для метода Рунге--Кутты:
\begin{align*}
 k_1&=-0{,}1, & k_2&=-0{,}095,\\
 k_3&=-0{,}09525, & k_4&=-0{,}090475.
\end{align*}
Отсюда
\begin{equation}
 y_1=1+\frac{-0{,}1+2(-0{,}095)+2(-0{,}09525)-0{,}090475}{6}
 =0{,}9048375.
\end{equation}
Точное значение в этой точке равно $e^{-0{,}1}=0{,}904837418$, что подтверждает высокую точность метода.

\section{Результаты вычислений}
В таблице приведены значения в каждом пятом узле сетки.

\begin{table}[H]
\centering
\caption{Сравнение численных решений}
\begin{tabular}{c|c|c|c|c}
\toprule
$x$ & Точное решение & Эйлер & Рунге--Кутта 4 & Милн \\
\midrule
0{,}0 & 1{,}00000000 & 1{,}00000000 & 1{,}00000000 & 1{,}00000000 \\
0{,}5 & 0{,}60653066 & 0{,}59049000 & 0{,}60653093 & 0{,}60653071 \\
1{,}0 & 0{,}36787944 & 0{,}34867844 & 0{,}36787977 & 0{,}36787909 \\
1{,}5 & 0{,}22313016 & 0{,}20589113 & 0{,}22313046 & 0{,}22312999 \\
2{,}0 & 0{,}13533528 & 0{,}12157665 & 0{,}13533553 & 0{,}13533490 \\
\bottomrule
\end{tabular}
\end{table}

\begin{table}[H]
\centering
\caption{Максимальные погрешности при $h=0{,}1$}
\begin{tabular}{l|c|c|c}
\toprule
Метод & Порядок & $\varepsilon_{\max}$ & Оценка Рунге \\
\midrule
Эйлер & 1 & $1{,}920100\cdot10^{-2}$ & $6{,}935502\cdot10^{-3}$ \\
Рунге--Кутта 4 & 4 & $3{,}332411\cdot10^{-7}$ & $1{,}536584\cdot10^{-8}$ \\
Милн & 4 & $3{,}981697\cdot10^{-7}$ & --- \\
\bottomrule
\end{tabular}
\end{table}

\section{Описание программы}
Программа имеет графический интерфейс на Streamlit. Пользователь выбирает одно из доступных ОДУ и задаёт границы интервала, шаг и требуемую точность. Для выбранной задачи последовательно вызываются функции методов Эйлера, Рунге--Кутты и Милна. Результаты отображаются в таблице вместе с точным решением, абсолютной погрешностью и оценкой по правилу Рунге. Дополнительно строятся графики численных и точного решений.

Основные вычислительные процедуры вынесены в отдельный модуль методов, структуры входных данных и результатов описаны с помощью классов данных, а расчёт оценок погрешности выполняется отдельными вспомогательными функциями. Такое разделение упрощает проверку формул и добавление новых задач Коши.

\section*{Вывод}
В работе реализованы и исследованы три численных метода решения задачи Коши. Метод Эйлера имеет первый порядок точности и при выбранном шаге заметно отклоняется от точного решения. Методы Рунге--Кутты четвёртого порядка и Милна дают значения, близкие к $e^{-x}$, а их максимальные погрешности имеют порядок $10^{-7}$. Следовательно, для получения высокой точности при умеренном числе узлов предпочтительнее использовать метод Рунге--Кутты или многошаговый метод Милна.

\end{document}